
\chapter{M002 --- Primitive Surface Construction}

\section{Stage 2 --- Mathematical Analysis}

\subsection{Objective}

Establish that the transformation from a Surface Frame to a Primitive
Surface is mathematically well defined independently of the CALM
implementation.

\subsection{Definition}

A Primitive Surface is the minimal two-dimensional periodic crystal
associated with a Surface Frame that preserves the crystallographic
orientation and atomic motif.

\subsection{Proposition M002.1}

Given a valid Surface Frame, there exists a corresponding primitive
surface lattice.

\subsection{Proof Strategy}

The proof will establish:

\begin{enumerate}
\item The adapted in-plane basis spans a rank-two lattice.
\item The projected crystal motif is periodic on that lattice.
\item A minimal periodic representative exists.
\item Equivalent primitive representatives differ only by gauge.
\end{enumerate}

\subsection{Open Mathematical Questions}

The following questions remain intentionally unresolved until the proof
is completed:

\begin{itemize}
\item Uniqueness of the primitive representative.
\item Interaction with crystallographic symmetry.
\item Treatment of equivalent surface terminations.
\end{itemize}

\subsection{Completion Criteria}

Stage 2 will be complete when the propositions above have been proved
without reference to the implementation.
